\textbf{Phần 1. Trắc nghiệm lựa chọn}
\setcounterpageref{ex}{0}

\begin{ex} %[5H3V18]
Một khu bảo tồn động vật hoang dã có diện tích $25\text{ km}^2$. Người ta ngăn ra $250\text{ ha}$ để làm khu chuồng nuôi bán tự nhiên, phần còn lại là rừng tự nhiên. Diện tích rừng tự nhiên là bao nhiêu héc-ta?
\choice
{$2\,000\text{ ha}$}
{\True $2\,250\text{ ha}$}
{$2\,475\text{ ha}$}
{$225\text{ ha}$}
\loigiai{
Đổi $25\text{ km}^2 = 2\,500\text{ ha}$.\\
Diện tích rừng tự nhiên là:
\[
2\,500 - 250 = 2\,250\text{ (ha)}.
\]
\textbf{Đáp án của câu hỏi là B}
}
\end{ex}

\begin{ex} %[5H3V18]
Một cánh đồng hình chữ nhật có chu vi $3\text{ km}$, chiều rộng bằng $\dfrac{2}{3}$ chiều dài. Diện tích cánh đồng đó tính theo héc-ta là
\choice
{$54\text{ ha}$}
{\True $54\text{ ha}$}
{$540\text{ ha}$}
{$216\text{ ha}$}
\loigiai{
Nửa chu vi cánh đồng là:
\[
3 : 2 = 1{,}5\text{ (km)} = 1\,500\text{ m}.
\]
Tổng số phần bằng nhau: $2 + 3 = 5$ phần.\\
Chiều dài cánh đồng: $1\,500 : 5 \times 3 = 900\text{ (m)}$.\\
Chiều rộng cánh đồng: $1\,500 - 900 = 600\text{ (m)}$.\\
Diện tích cánh đồng:
\[
900 \times 600 = 540\,000\text{ (m}^2\text{)} = 54\text{ ha}.
\]
\textbf{Đáp án của câu hỏi là A}
}
\end{ex}

\begin{ex} %[5H3V18]
Một mảnh đất hình bình hành có độ dài đáy là $180\text{ m}$, chiều cao bằng $\dfrac{5}{6}$ độ dài đáy. Diện tích mảnh đất đó bằng bao nhiêu héc-ta?
\choice
{$2{,}7\text{ ha}$}
{\True $2{,}7\text{ ha}$}
{$27\text{ ha}$}
{$0{,}27\text{ ha}$}
\loigiai{
Chiều cao mảnh đất là:
\[
180 \times \dfrac{5}{6} = 150\text{ (m)}.
\]
Diện tích mảnh đất là:
\[
180 \times 150 = 27\,000\text{ (m}^2\text{)}.
\]
Đổi sang héc-ta: $27\,000\text{ m}^2 = 2{,}7\text{ ha}$.\\
\textbf{Đáp án của câu hỏi là A}
}
\end{ex}

\begin{ex} %[5H3V18]
Người ta lát sàn một hội trường hình chữ nhật có kích thước $18\text{ m} \times 12\text{ m}$ bằng các viên gạch men hình vuông cạnh $60\text{ cm}$. Số viên gạch men cần dùng là
\choice
{$500$ viên}
{\True $600$ viên}
{$700$ viên}
{$800$ viên}
\loigiai{
Diện tích hội trường: $18 \times 12 = 216\text{ (m}^2\text{)} = 2\,160\,000\text{ cm}^2$.\\
Diện tích một viên gạch: $60 \times 60 = 3\,600\text{ (cm}^2\text{)}$.\\
Số viên gạch cần dùng là:
\[
2\,160\,000 : 3\,600 = 600\text{ (viên)}.
\]
\textbf{Đáp án của câu hỏi là B}
}
\end{ex}

\begin{ex} %[5H3V18]
Trong các số đo sau, số đo nào lớn nhất?
\choice
{$4\text{ ha } 50\text{ m}^2$}
{$4{,}5\text{ ha}$}
{\True $4\text{ ha } 500\text{ m}^2$}
{$40\,500\text{ m}^2$}
\loigiai{
Đổi tất cả về $\text{m}^2$:
\begin{itemize}
    \item $4\text{ ha } 50\text{ m}^2 = 40\,050\text{ m}^2$.
    \item $4{,}5\text{ ha} = 45\,000\text{ m}^2$.
    \item $4\text{ ha } 500\text{ m}^2 = 40\,500\text{ m}^2$.
    Khoan! $4{,}5\text{ ha} = 45\,000\text{ m}^2$ lớn hơn $40\,500\text{ m}^2$!
    Sửa lại: $4{,}5\text{ ha}$ lớn nhất!
\end{itemize}
Vì $45\,000\text{ m}^2 > 40\,500\text{ m}^2 > 40\,050\text{ m}^2$ nên $4{,}5\text{ ha}$ lớn nhất.\\
\textbf{Đáp án của câu hỏi là B}
}
\end{ex}

\begin{ex} %[5H3V18]
Một thửa ruộng hình chữ nhật có diện tích $1{,}2\text{ ha}$. Người ta chia thửa ruộng làm $3$ phần bằng nhau để trồng các giống lúa khác nhau. Diện tích mỗi phần là bao nhiêu mét vuông?
\choice
{$400\text{ m}^2$}
{\True $4\,000\text{ m}^2$}
{$40\,000\text{ m}^2$}
{$4\text{ m}^2$}
\loigiai{
Đổi $1{,}2\text{ ha} = 12\,000\text{ m}^2$.\\
Diện tích mỗi phần là:
\[
12\,000 : 3 = 4\,000\text{ (m}^2\text{)}.
\]
\textbf{Đáp án của câu hỏi là B}
}
\end{ex}

\begin{ex} %[5H3V18]
Một tấm tôn phẳng có diện tích $3\text{ m}^2$. Người ta gò một chiếc thùng không nắp có dạng hình hộp chữ nhật hết $240\text{ dm}^2$ tôn. Phần tôn còn thừa là bao nhiêu đề-xi-mét vuông?
\choice
{$6\text{ dm}^2$}
{\True $60\text{ dm}^2$}
{$600\text{ dm}^2$}
{$160\text{ dm}^2$}
\loigiai{
Đổi $3\text{ m}^2 = 300\text{ dm}^2$.\\
Phần tôn còn thừa là:
\[
300 - 240 = 60\text{ (dm}^2\text{)}.
\]
\textbf{Đáp án của câu hỏi là B}
}
\end{ex}

\begin{ex} %[5H3V18]
Một khu rừng cao su hình chữ nhật có chiều dài $3\text{ km}$, chiều rộng bằng $\dfrac{1}{3}$ chiều dài. Trung bình cứ mỗi héc-ta rừng cho thu hoạch $2{,}5\text{ tấn}$ mủ cao su mỗi năm. Hằng năm khu rừng đó cho sản lượng mủ cao su là
\choice
{$250\text{ tấn}$}
{\True $750\text{ tấn}$}
{$500\text{ tấn}$}
{$1\,000\text{ tấn}$}
\loigiai{
Chiều rộng khu rừng là:
\[
3 \times \dfrac{1}{3} = 1\text{ (km)}.
\]
Diện tích khu rừng là:
\[
3 \times 1 = 3\text{ (km}^2\text{)} = 300\text{ ha}.
\]
Sản lượng mủ cao su hằng năm là:
\[
300 \times 2{,}5 = 750\text{ (tấn)}.
\]
\textbf{Đáp án của câu hỏi là B}
}
\end{ex}

\begin{ex} %[5H3V18]
Một hình thoi có diện tích bằng $4{,}5\text{ m}^2$, độ dài một đường chéo là $3\text{ m}$. Độ dài đường chéo còn lại là
\choice
{$1{,}5\text{ m}$}
{\True $3\text{ m}$}
{$4{,}5\text{ m}$}
{$6\text{ m}$}
\loigiai{
Độ dài đường chéo còn lại là:
\[
4{,}5 \times 2 : 3 = 3\text{ (m)}.
\]
\textbf{Đáp án của câu hỏi là B}
}
\end{ex}

\begin{ex} %[5H3V18]
Đổi số đo diện tích $0{,}06\text{ km}^2$ ra mét vuông ta được
\choice
{$600\text{ m}^2$}
{$6\,000\text{ m}^2$}
{\True $60\,000\text{ m}^2$}
{$600\,000\text{ m}^2$}
\loigiai{
Vì $1\text{ km}^2 = 1\,000\,000\text{ m}^2$ nên:
\[
0{,}06 \times 1\,000\,000 = 60\,000\text{ (m}^2\text{)}.
\]
\textbf{Đáp án của câu hỏi là C}
}
\end{ex}

\begin{ex} %[5H3V18]
Một công viên có diện tích $16\text{ ha}$. Ban quản lý dành $\dfrac{1}{4}$ diện tích làm hồ nước, $\dfrac{3}{8}$ diện tích trồng cây xanh, phần còn lại xây dựng khu vui chơi trẻ em và đường đi bộ. Diện tích khu vui chơi và đường đi bộ là
\choice
{$4\text{ ha}$}
{\True $6\text{ ha}$}
{$8\text{ ha}$}
{$10\text{ ha}$}
\loigiai{
Phân số chỉ diện tích hồ nước và cây xanh là:
\[
\dfrac{1}{4} + \dfrac{3}{8} = \dfrac{2}{8} + \dfrac{3}{8} = \dfrac{5}{8}\text{ (diện tích)}.
\]
Phân số chỉ diện tích khu vui chơi và đường đi bộ là:
\[
1 - \dfrac{5}{8} = \dfrac{3}{8}\text{ (diện tích)}.
\]
Diện tích khu vui chơi và đường đi bộ là:
\[
16 \times \dfrac{3}{8} = 6\text{ (ha)}.
\]
\textbf{Đáp án của câu hỏi là B}
}
\end{ex}

\begin{ex} %[5H3V18]
Số thích hợp điền vào dấu hỏi chấm để $2\text{ ha } 45\text{ m}^2 = ?\text{ m}^2$ là
\choice
{$245$}
{$2\,045$}
{\True $20\,045$}
{$200\,045$}
\loigiai{
Vì $1\text{ ha} = 10\,000\text{ m}^2$ nên:
\[
2\text{ ha } 45\text{ m}^2 = 20\,000\text{ m}^2 + 45\text{ m}^2 = 20\,045\text{ m}^2.
\]
\textbf{Đáp án của câu hỏi là C}
}
\end{ex}

\textbf{Phần 2. Trắc nghiệm đúng sai}
\setcounterpageref{ex}{0}

\begin{ex} %[5H3V18]
Một nông trường tổng kết diện tích đất canh tác của ba đội sản xuất:
- Đội 1: $1\text{ km}^2\ 20\text{ ha}$
- Đội 2: $95\text{ ha}$
- Đội 3: $850\,000\text{ m}^2$
Xét tính đúng sai của các khẳng định sau:
\choiceTF
{\True Diện tích đất của Đội 1 là $120\text{ ha}$.}
{\True Diện tích đất của Đội 3 là $85\text{ ha}$.}
{Đội 2 có diện tích đất canh tác lớn nhất trong ba đội.}
{\True Tổng diện tích đất canh tác của cả ba đội là $300\text{ ha} = 3\text{ km}^2$.}
\loigiai{
\begin{itemize}
    \item Đội 1: $1\text{ km}^2\ 20\text{ ha} = 100\text{ ha} + 20\text{ ha} = 120\text{ ha}$. Khẳng định a \textbf{Đúng}.
    \item Đội 3: $850\,000\text{ m}^2 = 85\text{ ha}$. Khẳng định b \textbf{Đúng}.
    \item So sánh: $85\text{ ha} < 95\text{ ha} < 120\text{ ha}$ nên Đội 1 lớn nhất, không phải Đội 2. Khẳng định c \textbf{Sai}.
    \item Tổng diện tích: $120 + 95 + 85 = 300\text{ ha} = 3\text{ km}^2$. Khẳng định d \textbf{Đúng}.
\end{itemize}
\begin{itemize}
    \item a) \textbf{Đúng}.
    \item b) \textbf{Đúng}.
    \item c) \textbf{Sai}.
    \item d) \textbf{Đúng}.
\end{itemize}
}
\end{ex}

\begin{ex} %[5H3V18]
Người ta quy hoạch một khu du lịch sinh thái rộng $20\text{ ha}$ như sau:
- Khu đồi thông chiếm $\dfrac{2}{5}$ diện tích.
- Khu hồ bơi và thể thao chiếm $\dfrac{1}{4}$ diện tích.
- Phần còn lại là khu cắm trại và nhà hàng.
Xét tính đúng sai của các khẳng định sau:
\choiceTF
{\True Diện tích khu đồi thông là $8\text{ ha}$.}
{\True Diện tích khu hồ bơi và thể thao là $5\text{ ha}$.}
{Diện tích khu cắm trại và nhà hàng là $8\text{ ha}$.}
{\True Khu đồi thông có diện tích gấp hơn $1{,}5$ lần diện tích khu hồ bơi.}
\loigiai{
\begin{itemize}
    \item Khu đồi thông: $20 \times \dfrac{2}{5} = 8\text{ (ha)}$. Khẳng định a \textbf{Đúng}.
    \item Khu thể thao: $20 \times \dfrac{1}{4} = 5\text{ (ha)}$. Khẳng định b \textbf{Đúng}.
    \item Khu cắm trại và nhà hàng: $20 - (8 + 5) = 7\text{ (ha)} 
e 8\text{ ha}$. Khẳng định c \textbf{Sai}.
    \item Tỉ số diện tích: $8 : 5 = 1{,}6 > 1{,}5$. Khẳng định d \textbf{Đúng}.
\end{itemize}
\begin{itemize}
    \item a) \textbf{Đúng}.
    \item b) \textbf{Đúng}.
    \item c) \textbf{Sai}: Diện tích là $7\text{ ha}$.
    \item d) \textbf{Đúng}.
\end{itemize}
}
\end{ex}

\textbf{Phần 3. Trả lời ngắn}
\setcounterpageref{ex}{0}

\begin{ex} %[5H3V18]
Một cánh đồng hình chữ nhật có chiều dài $1\,200\text{ m}$, chiều rộng $500\text{ m}$. Hãy tính diện tích cánh đồng đó theo đơn vị héc-ta.
\shortans{60}
\loigiai{
\textbf{Đáp số:} $60$\\
Diện tích cánh đồng là:
\[
1\,200 \times 500 = 600\,000\text{ (m}^2\text{)}.
\]
Đổi sang héc-ta: $600\,000\text{ m}^2 = 60\text{ ha}$.
}
\end{ex}

\begin{ex} %[5H3V18]
Đổi $3\text{ km}^2\ 45\text{ ha}$ ra héc-ta. Kết quả là bao nhiêu héc-ta?
\shortans{345}
\loigiai{
\textbf{Đáp số:} $345$\\
Vì $1\text{ km}^2 = 100\text{ ha}$ nên:
\[
3\text{ km}^2\ 45\text{ ha} = 300\text{ ha} + 45\text{ ha} = 345\text{ ha}.
\]
}
\end{ex}

\begin{ex} %[5H3V18]
Một hội trường diện tích $280\text{ m}^2$ được lát bằng các viên gạch vuông cạnh $40\text{ cm}$. Biết mỗi hộp gạch có $5$ viên. Hỏi cần mua bao nhiêu hộp gạch như thế?
\shortans{350}
\loigiai{
\textbf{Đáp số:} $350$\\
Diện tích hội trường: $280\text{ m}^2 = 2\,800\,000\text{ cm}^2$.\\
Diện tích một viên gạch: $40 \times 40 = 1\,600\text{ cm}^2$.\\
Số viên gạch cần dùng: $2\,800\,000 : 1\,600 = 1\,750$ (viên).\\
Số hộp gạch cần mua: $1\,750 : 5 = 350$ (hộp).
}
\end{ex}

\begin{ex} %[5H3V18]
Một vườn cây hình vuông có chu vi $1\,200\text{ m}$. Diện tích vườn cây đó bằng bao nhiêu héc-ta?
\shortans{9}
\loigiai{
\textbf{Đáp số:} $9$\\
Cạnh vườn cây hình vuông: $1\,200 : 4 = 300\text{ (m)}$.\\
Diện tích vườn cây: $300 \times 300 = 90\,000\text{ (m}^2\text{)} = 9\text{ ha}$.
}
\end{ex}

\begin{ex} %[5H3V18]
Đổi số đo diện tích $0{,}75\text{ km}^2$ ra héc-ta. Kết quả là bao nhiêu héc-ta?
\shortans{75}
\loigiai{
\textbf{Đáp số:} $75$\\
Vì $1\text{ km}^2 = 100\text{ ha}$ nên:
\[
0{,}75\text{ km}^2 = 0{,}75 \times 100 = 75\text{ ha}.
\]
}
\end{ex}

\begin{ex} %[5H3V18]
Một mảnh đất hình thoi có diện tích $450\text{ m}^2$. Độ dài một đường chéo là $30\text{ m}$. Độ dài đường chéo còn lại là bao nhiêu mét?
\shortans{30}
\loigiai{
\textbf{Đáp số:} $30$\\
Độ dài đường chéo còn lại là:
\[
450 \times 2 : 30 = 30\text{ (m)}.
\]
}
\end{ex}
